\begin{table}[!htbp]
\centering
\footnotesize
\begin{threeparttable}
\caption{Empirical avalanche-size distribution. An avalanche of size $s$ is a calendar month in which $s$ countries simultaneously exceed their country-specific stress threshold ($\sigma=1.5$).}
\label{tab:avalanche_distribution}
\begin{tabular}{rrrr}
\toprule
$s$ & Frequency & Probability $P(S=s)$ & CCDF $P(S \geq s)$ \\
\midrule
1 & 18 & 0.450 & 1.000 \\
2 & 6 & 0.150 & 0.550 \\
3 & 2 & 0.050 & 0.400 \\
4 & 2 & 0.050 & 0.350 \\
5 & 2 & 0.050 & 0.300 \\
6 & 3 & 0.075 & 0.250 \\
8 & 2 & 0.050 & 0.175 \\
9 & 1 & 0.025 & 0.125 \\
10 & 1 & 0.025 & 0.100 \\
11 & 3 & 0.075 & 0.075 \\
\bottomrule
\end{tabular}
\begin{tablenotes}\footnotesize
\item Total avalanche months: $T_a=40$. Total country-stress events: $\sum_s s\cdot n_s=1540$.
\item The distribution is heavy-tailed with mass at the extremes: $P(S\geq 6)=0.250$, $P(S=11)=0.075$.
\item All $S=11$ events correspond to global systemic shocks (post-Lehman 2008, COVID-19 2020).
\end{tablenotes}
\end{threeparttable}
\end{table}
